%––––––––––––––––––––––––––––––––––––––––
%
%
%––––––––––––––––––––––––––––––––––––––––
\pagebreak
\unnumberedSection{BK - Remove simulator breakpoint}

%––––––––––––––––––––––––––––––––––––––––
\begin{iPageDescEntry}{Syntax}
\texttt{bk  <bNum> [ , <modNum> ] }
\end{iPageDescEntry}

%––––––––––––––––––––––––––––––––––––––––
\begin{iPageDescEntry}{Description}

The \texttt{BK} command removes a simulator breakpoint.

The \texttt{bNum} parameter specifies the breakpoint number as shown by
the \texttt{BL} command. The optional \texttt{modNum} parameter specifies the
module from which the breakpoint is to be removed. A module number of
\texttt{-1} removes the breakpoint from all modules.

If the breakpoint is associated with more than one module, removing it
from one module does not remove the breakpoint itself. The module is
instead removed from the breakpoint’s module mask. The breakpoint is
removed from the breakpoint table only when no modules remain
associated with it.

\end{iPageDescEntry}

%––––––––––––––––––––––––––––––––––––––––
\begin{iPageDescEntry}{Example}

The following examples demonstrate the \texttt{BK} command.

\begin{lstlisting}[style=iPageOpStyle]
(6) ->bl
 Idx  Type   State   Adr                  Modules         
 0    CODE   E       0x00_0000_1000       0000_0000_0000_0010
 1    CODE   E       0x00_0000_2000       0000_0000_0000_0010
 2    DATA   E       0x00_0000_3000:0004  0000_0000_0000_0010
 3    DATA   E       0x00_0000_3100:0100  0000_0000_0000_0020
 (7) ->bk 1
 (8) ->bl
 Idx  Type   State   Adr                  Modules         
 0    CODE   E       0x00_0000_1000       0000_0000_0000_0010
 2    DATA   E       0x00_0000_3000:0004  0000_0000_0000_0010
 3    DATA   E       0x00_0000_3100:0100  0000_0000_0000_0020
 (9) ->
\end{lstlisting}

\end{iPageDescEntry}

%––––––––––––––––––––––––––––––––––––––––
\begin{iPageDescEntry}{Notes}

The breakpoint number of removed breakpoint is re-used for the next breakpoint create.

\end{iPageDescEntry}